To prove the reverse inclusions, let us henceforth suppose
\(\mathfrak{g}=\Lie(G/S,\mathcal{M})\) commutative; then \(\mathfrak{g}^{K}\) and
\((\mathfrak{g}/\mathfrak{k})^{K}\) are commutative \(S\)-groups. Let
\(X\in\mathfrak{g}(S)\) and \(X\) its image in
\((\mathfrak{g}/\mathfrak{k})(S)\), suppose that
\(X\in(\mathfrak{g}/\mathfrak{k})^{K}(S)\) (\resp~\(X\in\mathfrak{g}^{K}(S)\))
and let us show that \(X\in\mathfrak{n}(S)\)
(\resp~\(X\in\mathfrak{z}(S)\)).

Let \(f\colon S'\to I_S(\mathcal{M})\); let us show that the preceding
condition~\((*)\) is satisfied for every \(u\in K(S')\). Consider the
cartesian square
\[
\xymatrix{
I_{S'}(\mathcal{M}) \ar[r]^{p} \ar[d]_{\rho'}
  & I_S(\mathcal{M}) \ar[d]^{\rho} \\
S' \ar[r]_{\rho\circ f} & S
}
\]
and let \(h\) be the section of \(\rho'\) defined by~\(f\). It suffices to
show that, for every \(v\in K(I_{S'}(\mathcal{M}))\), one has
\[
\left.
\begin{aligned}
p^{*}(X^{\alpha})\cdot v\cdot p^{*}(X^{-\alpha})\cdot v^{-1}
  &\in K\bigl(I_{S'}(\mathcal{M})\bigr)\\
\text{resp.}\quad
p^{*}(X)\cdot v\cdot p^{*}(X^{-1})\cdot v^{-1}
  &=1
\end{aligned}
\right\}
\quad(**).
\]
Indeed, taking \(v=\rho^{\prime*}(u)\) and applying \(h^{*}\) to~\((**)\),
one obtains~\((*)\), since \(\rho'\circ h=\mathrm{id}_{S'}\) and
\(p\circ h=f\).

Let us now show~\((**)\); to simplify, one will write \(X\) in place of
\(p^{*}(X)\). Every \(v\in K(I_{S'}(\mathcal{M}))\) may be written uniquely
\(Y\cdot k\) where \(Y\in\Lie(K/S,\mathcal{M})(S')\) and \(k\in K(S')\).
The expression \(X^{\alpha}\cdot u\cdot X^{-\alpha}\cdot u^{-1}\) then
% typo?: kept as printed (u, not the running variable v)
becomes \(X^{\alpha}\cdot Y\cdot(k\cdot X^{-\alpha}\cdot k^{-1})\cdot Y^{-1}\)
which, as \(\Lie(G/S,\mathcal{M})\) is commutative, is written
\(X^{\alpha}\,\Ad(k)(X^{-\alpha})\). Now the latter lies a priori in
\(\Lie(G/S,\mathcal{M})(S')\); taking Lemma~5.2.3.0 into account, the
condition~\((**)\) for \(X\) therefore becomes: for every
\(k\in K(S')\),
\[
\begin{cases}
\Ad(k)(X)=X
  & \text{if }H=Z;\\
X^{\alpha}\,\Ad(k)(X^{-\alpha})\in\Lie(K/S,\mathcal{M})(S')
  & \text{if }H=N.
\end{cases}
\]
When \(H=Z\), this condition is indeed a consequence of the hypothesis
\(X\in\mathfrak{g}^{K}(S)\).
When \(H=N\), the condition may also be written:
\[
(*')\qquad
\Ad(k)(X)=X
\qquad\text{and}\qquad
\Ad(k)(X^{-1})=X^{-1},
\]
now the second condition of~\((*')\) is a consequence of the first, since
the operation of \(K\) on \(\mathfrak{g}/\mathfrak{k}\) respects the group
structure of the latter. Thus, when \(H=N\), the condition~\((**)\) for
\(X\) is indeed a consequence of the hypothesis
\(X\in(\mathfrak{g}/\mathfrak{k})^{K}(S)\). This proves~(i) and~(ii).

To prove~(iii)--(v), let us write
\(\mathfrak{g}(\mathcal{M})=\Lie(G/S,\mathcal{M})\) and let us define
likewise \(\mathfrak{k}(\mathcal{M})\), \(\mathfrak{z}(\mathcal{M})\) and
\(\mathfrak{n}(\mathcal{M})\). If \(G/S\) satisfies~(E), then
\(\mathfrak{g}(\mathcal{M}\oplus\mathcal{N})
\simeq\mathfrak{g}(\mathcal{M})\times_S\mathfrak{g}(\mathcal{N})\)
and hence, by~5.2.2~(i),
\[
\mathfrak{g}(\mathcal{M}\oplus\mathcal{N})^{K}
\simeq
\mathfrak{g}(\mathcal{M})^{K}\times_S\mathfrak{g}(\mathcal{N})^{K}
\]
whence
\(\mathfrak{z}(\mathcal{M}\oplus\mathcal{N})
\simeq\mathfrak{z}(\mathcal{M})\times_S\mathfrak{z}(\mathcal{S})\),
% typo?: kept as printed \(\mathfrak{z}(\mathcal{S})\); the parallel isomorphism has \(\mathcal{N}\)
hence \(Z\) satisfies~(E). If moreover \(K/S\) satisfies~(E), one obtains
successively isomorphisms:
\begin{align*}
(\mathfrak{g}/\mathfrak{k})(\mathcal{M}\oplus\mathcal{N})
  &\simeq
  (\mathfrak{g}/\mathfrak{k})(\mathcal{M})
  \times_S
  (\mathfrak{g}/\mathfrak{k})(\mathcal{N})\\
(\mathfrak{g}/\mathfrak{k})^{K}(\mathcal{M}\oplus\mathcal{N})
  &\simeq
  (\mathfrak{g}/\mathfrak{k})(\mathcal{M})^{K}
  \times_S
  (\mathfrak{g}/\mathfrak{k})(\mathcal{N})^{K},
\end{align*}
then a commutative diagram:
\[
\xymatrix@C=1.35em@R=1.55em{
0 \ar[r]
  & \mathfrak{k}(\mathcal{M}\oplus\mathcal{N}) \ar[r] \ar[d]^{\wr}
  & \mathfrak{n}(\mathcal{M}\oplus\mathcal{N}) \ar[r] \ar[d]
  & (\mathfrak{n}/\mathfrak{k})(\mathcal{M}\oplus\mathcal{N}) \ar[r] \ar[d]^{\wr}
  & 0
\\
0 \ar[r]
  & \mathfrak{k}(\mathcal{M})\times_S\mathfrak{k}(\mathcal{N}) \ar[r]
  & \mathfrak{n}(\mathcal{M})\times_S\mathfrak{n}(\mathcal{N}) \ar[r]
  & (\mathfrak{n}/\mathfrak{k})(\mathcal{M})
      \times_S(\mathfrak{n}/\mathfrak{k})(\mathcal{N}) \ar[r]
  & 0
}
\]
whence it follows that
\(\mathfrak{n}(\mathcal{M}\oplus\mathcal{N})
\simeq\mathfrak{n}(\mathcal{M})\times_S\mathfrak{n}(\mathcal{N})\),
hence \(N\) satisfies~(E).

Henceforth, let us write \(\mathfrak{g}=\mathfrak{g}(\OS)=\Lie(G/S)\) and
let us define likewise \(\mathfrak{k}\), \(\mathfrak{z}\), \(\mathfrak{n}\).
If \(G\) is \emph{good}, \(\mathfrak{g}\) is a \emph{good} \(\OS\)-module,
hence, by~5.2.2~(ii), so is \(\mathfrak{z}\simeq\mathfrak{g}^{K}\), hence
\(Z\) (which satisfies~(E) by~(iii)) is \emph{good}. If moreover \(K\) is
\emph{good}, then \(\mathfrak{k}\) is a \emph{good} \(\OS\)-module, hence,
by~5.2.1~(iv) and~5.2.2~(ii), the same is true of
\(\mathfrak{g}/\mathfrak{k}\) and \((\mathfrak{g}/\mathfrak{k})^{K}\). In
view of the exact sequence
\[
0\longrightarrow\mathfrak{k}\longrightarrow\mathfrak{n}
\longrightarrow(\mathfrak{g}/\mathfrak{k})^{K}\longrightarrow 0,
\]
one obtains, by~5.2.1~(iv) again, that \(\mathfrak{n}\) is \emph{good}.
Finally, if in addition to the preceding conditions \(G\) is \emph{very
good}, that is, if one has identically \([x,x]=0\) for every
\(x\in\mathfrak{g}\), it is clear that \(Z\) and \(N\) are \emph{very
good}. This proves~(iii), (iv) and~(v).

\begin{corollary}{5.2.3.1}
\label{II.5.2.3.1}
One has \(\Lie(\Centr(G)/S)=\Lie(G/S)^{G}\) if the group law of
\(\Lie(G/S)\) is commutative.
\end{corollary}

\begin{corollary}{5.2.3.2}
\label{II.5.2.3.2}
If the group law of \(\Lie(G/S)\) is commutative and if \(K\) is an
invariant subgroup of \(G\), then
\[
\bigl(\Lie(G/S)\big/\Lie(K/S)\bigr)^{K}
=\Lie(G/S)\big/\Lie(K/S).
\]
\end{corollary}

\par\addvspace{0.55\baselineskip}%
\noindent\textbf{5.3. Linear representations.}\ ---
\label{II.5.3}
Let \(G\) be a \emph{good} \(S\)-group operating linearly on a \emph{good}
\(\OS\)-module \(F\) via
\[
\rho\colon G\longrightarrow\Aut_{\OS\text{-mod.}}(F).
\]
One has defined (4.7.1 and~4.8.2) a corresponding linear representation
\[
\rho'\colon\Lie(G/S)\longrightarrow\End_{\OS\text{-mod.}}(F).
\]
The sub-\(S\)-groups \(\Norm_G(E)\) and \(\Centr_G(E)\) are defined for every
subset \(E\) of \(F\), for example for every sub-\(\OS\)-module \(E\) of~\(F\).

\begin{definition}{5.3.0}
\label{II.5.3.0}
One will set, analogously: for every \(S'\to S\),
\begin{align*}
\Norm_{\Lie(G/S)}(E)(S')
&=
\bigl\{X\in\Lie(G/S)(S')\mid \rho'(X)E_{S'}\subset E_{S'}\bigr\};\\
\Centr_{\Lie(G/S)}(E)(S')
&=
\bigl\{X\in\Lie(G/S)(S')\mid \rho'(X)E_{S'}=0\bigr\}.
\end{align*}
(Let us note that this construction is carried out for every linear
representation of an \(\OS\)-``Lie algebra'' (in the sense of~4.9) and that
the two subobjects constructed are sub-\(\OS\)-modules stable under the
bracket).
\end{definition}

\begin{theorem}{5.3.1}
\label{II.5.3.1}
Let \(G\) be a good \(S\)-group operating linearly on a good \(\OS\)-module
\(F\) and let \(E\) be a sub-\(\OS\)-module of~\(F\).
\begin{enumeratei}
\item One has \(\Lie(\Centr_G(E)/S)=\Centr_{\Lie(G/S)}(E)\) and
\(\Centr_G(E)\) is a good \(S\)-group; it is very good if \(G\)~is.
\item Suppose that \(E\) is a good \(\OS\)-module. Then%
{\renewcommand{\thefootnote}{(98)}\nde{One has corrected the original,
which stated the following equality without hypotheses on~\(E\); see~6.5.1
for counterexamples.}}
\(\Lie(\Norm_G(E)/S)=\Norm_{\Lie(G/S)}(E)\) and \(\Norm_G(E)\) is a good
\(S\)-group; it is very good if \(G\)~is.
\end{enumeratei}
\end{theorem}

The proof is left to the reader.

\numpara{5.3.2}
\label{II.5.3.2}
Let \(G\) be a good \(S\)-group; the foregoing applies in particular to the
case where one takes for \(\rho\) the adjoint representation of~\(G\). Let
\(E\) be a \emph{good}%
{\renewcommand{\thefootnote}{(99)}\nde{One has added this hypothesis;
\cf~6.5.1.}}
submodule of \(\Lie(G/S)\); one therefore associates with it two subgroups
of~\(G\), its centralizer and its normalizer. By~5.3.1, their Lie algebras
are respectively the centralizer and the normalizer of \(E\) in
\(\Lie(G/S)\), computed as usual by means of the bracket:
\begin{align*}
\Centr_{\Lie(G/S)}(E)(S')
&=
\bigl\{X\in\Lie(G/S)(S')\mid [X,E_{S'}]=0\bigr\}.\\
\Norm_{\Lie(G/S)}(E)(S')
&=
\bigl\{X\in\Lie(G/S)(S')\mid [X,E_{S'}]\subset E_{S'}\bigr\}.
\end{align*}

\numpara{5.3.3}
\label{II.5.3.3}
Let \(K\) be a sub-\(S\)-group of~\(G\), then \(\Lie(K/S)\) is a
sub-\(\OS\)-module of \(\Lie(G/S)\); suppose that \(\Lie(K/S)\) is a
\emph{good} \(\OS\)-module%
{\renewcommand{\thefootnote}{(99)}\ndemark}
(which is the case if \(K\) is a \emph{good} \(S\)-group). One evidently has
\begin{align*}
\Centr_G(K)&\subset\Centr_G\bigl(\Lie(K/S)\bigr)\\
\Norm_G(K)&\subset\Norm_G\bigl(\Lie(K/S)\bigr),
\end{align*}
whence, by~5.3.1,
\begin{align*}
\Lie(\Centr_G(K)/S)
  &\subset\Centr_{\Lie(G/S)}\bigl(\Lie(K/S)\bigr)\\
\Lie(\Norm_G(K)/S)
  &\subset\Norm_{\Lie(G/S)}\bigl(\Lie(K/S)\bigr),
\end{align*}
but none of these four inclusions is a priori an identity; we shall see
many examples of this later.

It follows in particular from these inclusions that if \(K\) is an
\emph{invariant subgroup} of~\(G\), then \(\Lie(K/S)\) is an \emph{ideal}
of \(\Lie(G/S)\).

\sgasection{6}{Miscellaneous remarks}
\label{II.sec.6}

\numpara{6.1}
\label{II.6.1}
One can define the bracket of two infinitesimal automorphisms for an
\(S\)-functor \(X\) which is not necessarily a group. It suffices to apply
the results of this exposé to the group \(\Aut_S(X)\). In order to arrive
at a pleasant formalism, one is led to suppose \(X\) \emph{good}, that is,
to suppose that the \(\mathcal{O}_X\)-module \(\mathrm{T}_{X/S}\) is
\emph{good} (if \(X\) is an \(S\)-group, this definition evidently coincides
with definition~4.6).

\numpara{6.2}
\label{II.6.2}
There exist functors possessing infinitesimal endomorphisms which are not
automorphisms, hence \emph{a fortiori} not satisfying
condition~(E).%
{\renewcommand{\thefootnote}{(100)}\nde{One has set out the original in
more detail in what follows. In particular, the correct definition of the
``free abelian monoid on a pointed set'' was pointed out to us by
O.~Gabber.}}
For every pointed set \((E,x_0)\), let \(\mathrm{MA}(E)\) be the free
abelian monoid generated by \(E\) and let \(\mathrm{MA}_{\mathrm{P}}(E,x_0)\)
be the abelian monoid obtained by quotienting \(\mathrm{MA}(E)\) by the
equivalence relation generated by the relation \(m\sim x_0+m\). Then
\((E,x_0)\mapsto\mathrm{MA}_{\mathrm{P}}(E,x_0)\) is the left adjoint of the
forgetful functor from the category of abelian monoids to that of pointed
sets; one will say that \(\mathrm{MA}_{\mathrm{P}}(E,x_0)\) is the ``free
abelian monoid on the pointed set \((E,x_0)\)''.

Let us then take for \(X\) the functor which to every scheme \(S\)
associates the free abelian monoid on the set \(\mathrm{O}(S)\), pointed
by the zero element. Each morphism
\(f\colon S\to I_{\ZZ}=\Spec(\ZZ[\varepsilon])\) corresponds to an element
\(u_f\) of square zero of \(\mathrm{O}(S)\), hence defines an endomorphism
of \(X(S)\) by \(x\mapsto x+u_f\) (sum in
\(\mathrm{MA}_{\mathrm{P}}(\mathrm{O}(S),0)\)). One thus obtains an
endomorphism \(\phi\) of \(X_{I_{\ZZ}}=X\times_{\ZZ}I_{\ZZ}\), defined as
follows. For every \(f\in I_{\ZZ}(S)\) and \(x\in X(S)\),
\[
\phi\bigl((x,f)\bigr)=(x+u_f,f).
\]
If \(f_0\colon S\to I_{\ZZ}\) is the composite of the structural morphism
\(S\to\Spec(\ZZ)\) and of the zero section of \(I_{\ZZ}\), the corresponding
element is \(u_{f_0}=0\), and hence
\(\phi\bigl((x,f_0)\bigr)=(x,f_0)\), since \(x+0=x\) in
\(\mathrm{MA}_{\mathrm{P}}(\mathrm{O}(S),0)\). This shows that \(\phi\)
induces the identity on~\(X\); it is therefore an infinitesimal
endomorphism of \(X\) which is evidently not an automorphism.

\numpara{6.3}
\label{II.6.3}
There exist modules which are not \emph{good}. On the one hand, one can
slightly modify the preceding counterexample:%
{\renewcommand{\thefootnote}{(101)}\nde{One has added what follows.}}
let us take for \(F(S)\) the free \(\mathrm{O}(S)\)-module with basis the
elements of \(\mathrm{O}(S)\); then \(F\) does not satisfy condition~(E)
with respect to \(\Spec(\ZZ)\); moreover, let \(G\) be the \(\ZZ\)-group
defined by \(G(S)=\GL_n(R(S))\), where \(n\geqslant 2\) and \(R(S)\) is
the \(\mathrm{O}(S)\)-algebra of the abelian group \(\mathrm{O}(S)\); then
the \(\ZZ\)-group \(\Lie(G/\ZZ)\) is not commutative.

{\renewcommand{\thefootnote}{(102)}\nde{In what follows, one has set out
example~a) in more detail and added example~b).}}%
On the other hand, one can give the following counterexamples. Let
\(S=\Spec(k)\), \(k\) a field of characteristic \(p>0\).

a)~Let \(F\) be the \(\OS\)-module which to every \(S\)-scheme \(T\)
associates \(F(T)=\Gamma(T,\mathcal{O}_T)\) equipped with the structure of
\(\mathrm{O}(T)\)-module obtained by making the scalars operate via the
\(p\)-th power, that is, \(r\cdot f=r^{p}f\), for \(r\in\mathrm{O}(T)\),
\(f\in F(T)\). As an \(S\)-group functor, \(F\) is isomorphic to
\(\mathbf{G}_{\mathrm{a},S}\). Hence \(F/S\) satisfies~(E) and
\(\Lie(F/S)\) is identified with
\(\Lie(\mathbf{G}_{\mathrm{a},S}/S)\simeq\OS\). Then, the canonical
morphism \(F\to\Lie(F/S)\) is,
for every \(T\to S\), the identity map \(F(T)\to\mathrm{O}(T)\): it does
respect the structure of abelian group but not the structure of module.
Hence \(F\) is not \emph{good} (\cf~4.4.1).

b)~Let \(\alpha_{p,k}\) be the \(k\)-group functor which to every
\(S\)-scheme \(T\) associates
\[
\alpha_{p,k}(T)=\{x\in\mathrm{O}(T)\mid x^{p}=0\},
\]
it is represented by \(\Spec(k[X]/(X^{p}))\) hence it is a \emph{very good}
\(S\)-group; it is also equipped with a structure of \(\OS\)-module, but it
is not a \emph{good} \(\OS\)-module, for the canonical morphism
\(\alpha_{p,k}\to\Lie(\alpha_{p,k}/k)=\mathbf{G}_{\mathrm{a},k}\) is not
bijective.

\numpara{6.4}
\label{II.6.4}
Let \(G\) be a group functor over~\(S\). One has by definition the following
implications
\[
\bigl(G/S\text{ satisfies~(E)}\bigr)
\Longleftarrow
\bigl(G\text{ is \emph{good}}\bigr)
\Longleftarrow
\bigl(G\text{ is \emph{very good}}\bigr).
\]
{\renewcommand{\thefootnote}{(103)}\nde{And \(G\) is very good if it is
representable (4.11). For criteria of representability, see for
example~[MO67]. On the other hand, for the automorphisms of an affine
algebraic group (over an algebraically closed field of characteristic~\(0\)),
let us cite~[HM69].}}
It would be interesting to prove or to give counterexamples for the
implications in the reverse direction.

\numpara{6.5}
\label{II.6.5}
{\renewcommand{\thefootnote}{(104)}\nde{One has added this paragraph.}}%
Let \(\mathrm{Nil}\) be the \(\ZZ\)-group functor defined as follows: for
every scheme \(S\), \(\mathrm{Nil}(S)\) is the ideal of \(\mathrm{O}(S)\)
consisting of the nilpotent elements, \ie
\[
\mathrm{Nil}(S)
=
\bigl\{x\in\mathrm{O}(S)\mid
\text{there exists }n\in\NN\text{ such that }x^{n}=0\bigr\}.
\]
(\(\mathrm{Nil}\) is \emph{very good} but is not representable). Let
\(\mathrm{Nil}^{2}\), \(\mathrm{O}_{\mathrm{red}}\) and \(F\) be the
\(\ZZ\)-group functors which to every scheme \(S\) associate, respectively,
the ideal \(\mathrm{Nil}(S)^{2}\) and
\[
\mathrm{O}_{\mathrm{red}}(S)=\mathrm{O}(S)/\mathrm{Nil}(S),
\qquad
F(S)=\mathrm{O}(S)/\mathrm{Nil}(S)^{2}.
\]
One sees easily that \(\Lie(\mathrm{O}_{\mathrm{red}}/\ZZ)=0\), hence the
\(\mathcal{O}_{\ZZ}\)-module \(\mathrm{O}_{\mathrm{red}}\) is not
\emph{good} (although it is a \emph{good} \(\ZZ\)-group). If
\(\mathcal{M}\), \(\mathcal{N}\) are free \(\ZZ\)-modules of finite rank,
one has
\[
\mathrm{Nil}^{2}\bigl(I_S(\mathcal{M}\oplus\mathcal{N})\bigr)
\simeq
\mathrm{Nil}^{2}(S)
\oplus\mathrm{Nil}(S)\otimes_{\ZZ}\mathcal{M}
\oplus\mathrm{Nil}(S)\otimes_{\ZZ}\mathcal{N}
\]
and hence
\[
F\bigl(I_S(\mathcal{M}\oplus\mathcal{N})\bigr)
\simeq
F(S)
\oplus\mathrm{O}_{\mathrm{red}}(S)\otimes_{\ZZ}\mathcal{M}
\oplus\mathrm{O}_{\mathrm{red}}(S)\otimes_{\ZZ}\mathcal{N}.
\]
One deduces from this, on the one hand, that the \(\ZZ\)-group functor \(F\)
satisfies condition~(E) and, on the other hand, that
\(\Lie(F/\ZZ)=\mathrm{O}_{\mathrm{red}}\); as the latter is not a
\emph{good} \(\mathcal{O}_{\ZZ}\)-module, this shows that \(F\) is an
example of a \(\ZZ\)-group which satisfies~(E) but which is not
\emph{good}.

\numpara{6.5.1}
\label{II.6.5.1}
Let us also give the counterexamples indicated in~5.3.1--5.3.3. Let \(S\)
be a scheme, \(F\) the \emph{good} \(\OS\)-module \(\OS^{\oplus 2}\)
equipped with the natural action of the \emph{good} \(S\)-group
\(G=\GL_{2,S}\), and \(E\) the sub-\(\OS\)-module of \(F\) consisting of the
pairs \((x_1,x_2)\) such that \(x_2\) is nilpotent. Set
\(N=\Norm_G(E)\). Then \(\Lie(N/S)=\Lie(G/S)\) whereas, for every
\(S'\to S\), one has
\[
\Norm_{\Lie(G/S)}(E)(S')
=
\left\{
\begin{pmatrix} a & b \\ x & c \end{pmatrix}
\;\middle|\;
a,\,b,\,c,\,x\in\mathrm{O}(S'),\ x\text{ nilpotent}
\right\}
\]
hence \(\Lie(\Norm_G(E)/S)\neq\Norm_{\Lie(G/S)}(E)\).

By considering the semidirect product \(G_1=F\cdot G\), one obtains an
analogous counterexample in which \(E\) is a sub-\(\OS\)-module of
\(\Lie(G_1/S)\); moreover, with the notations introduced above,
\(E=\Lie(K/S)\) where \(K\) is the subgroup
\(\OS\oplus\mathrm{Nil}^{2}\) of~\(F\) (that is, for every \(S'\to S\),
\(K(S')\) consists of the pairs \((x_1,x_2)\) such that
\(x_2\in\mathrm{Nil}(S')^{2}\)).

\makeatletter
\renewcommand{\@bibtitlestyle}{%
  \par\addvspace{1.1\baselineskip}%
  \noindent{\large\bfseries Bibliography}%
  {\renewcommand{\thefootnote}{(105)}\nde{Additional references cited in
  this Exposé.}}%
  \par\addvspace{0.45\baselineskip}%
}
\makeatother
\begin{thebibliography}{DG70}

\bibitem[DG70]{DG70}
M.~Demazure, P.~Gabriel, \emph{Groupes Algébriques}, Masson \&
North-Holland, 1970.

\bibitem[HM69]{HM69}
G.~Hochschild, G.~D.~Mostow, \emph{Automorphisms of affine algebraic
groups}, J.~Algebra \textbf{13} (1969), 535--543.

\bibitem[MO67]{MO67}
H.~Matsumura, F.~Oort, \emph{Representability of group functors and
automorphisms of algebraic schemes}, Invent.\ math.\ \textbf{4} (1967),
1--25.

\end{thebibliography}
